\documentclass[11pt]{article}
\usepackage[margin=1in]{geometry}
\usepackage{amsmath}
\usepackage{booktabs}
\usepackage{hyperref}
\usepackage{url}

\title{ArduPlane Elevon Mixing: Software Overview and Sign Convention}
\author{}
\date{}

\begin{document}
\maketitle

\section{Scope}

This note describes the fixed-wing Plane elevon mixer in the current ArduPilot
source. It covers the logical mixing performed in software; it does not prescribe
airframe-specific servo installation, linkage geometry, or flight tuning.
The principal implementation is
\url{ArduPlane/servos.cpp}, in \texttt{Plane::channel\_function\_mixer()} and
\texttt{Plane::servos\_output()}.

\section{Signal path in different modes}

Elevon mixing combines two logical control demands: aileron (roll) and elevator
(pitch). The source of those demands depends on the flight mode, but the
elevon output mixer is in the common servo-output path.

\begin{itemize}
  \item In \textbf{MANUAL}, \texttt{ModeManual::update()} writes the RC roll and
  pitch inputs, with the configured manual expo, to the logical aileron and
  elevator outputs. It does not run the normal attitude stabilization
  controllers for these axes.
  \item In \textbf{stabilized modes}, such as FBWA, the attitude controllers
  calculate the logical aileron and elevator demands. For example,
  \texttt{stabilize\_roll()} and \texttt{stabilize\_pitch()} write their
  controller outputs to those logical functions.
  \item In \textbf{automatic modes}, the mode's navigation and attitude control
  determine the demands; the same logical aileron/elevator output functions are
  ultimately used by the mixer. AUTO is therefore not a separate elevon mixing
  formula.
\end{itemize}

In the normal Plane output path, \texttt{servos\_output()} calls the mixer with
aileron and elevator as inputs and left and right elevons as outputs. Servo
function assignments route the resulting logical outputs to physical outputs:
\texttt{SERVO\textit{n}\_FUNCTION}=77 selects ElevonLeft and 78 selects
ElevonRight. The function values are defined in
\texttt{libraries/SRV\_Channel/SRV\_Channel.h}.

\section{Mixing equations}

Let \(a\) and \(e\) be the scaled aileron and elevator demands read by the
mixer. In the following, positive and negative refer to ArduPilot's logical
scaled values, not a transmitter's stick direction or a particular linkage.
The parameter \texttt{MIXING\_OFFSET} optionally adjusts the relative input
gains. For offset \(o\), in percentage points, define
\[
(a',e') =
\begin{cases}
(a,\;e(100-o)/100), & o < 0,\\
(a,\;e), & o = 0,\\
(a(100+o)/100,\;e), & o > 0.
\end{cases}
\]
Thus a positive offset increases the aileron contribution; a negative offset
increases the elevator contribution. The implementation then computes
\[
\begin{aligned}
L &= \operatorname{clip}_{[-4500,4500]}\bigl(g(e'-a')\bigr),\\
R &= \operatorname{clip}_{[-4500,4500]}\bigl(g(e'+a')\bigr),
\end{aligned}
\]
where \(L\) and \(R\) are the left and right scaled elevon outputs and \(g\)
is \texttt{MIXING\_GAIN}. The default gain is \(0.5\). Clipping is applied
separately to each output, so a sufficiently large combined demand can saturate
one or both elevons.

With zero offset and before clipping, the equations can also be inverted:
\[
e = \frac{L+R}{2g}, \qquad a = \frac{R-L}{2g}.
\]
This shows the symmetric component (pitch) in the sum and the differential
component (roll) in the right-minus-left difference.

\section{Deflection sign convention}

The mixer source comments that its output order is chosen so that
\textbf{non-reversed servos go ``up''}, and that the first output is the
\textbf{left} channel. Interpreting positive scaled elevon output as an upward
surface command for a correctly installed, non-reversed output gives:

\begin{center}
\begin{tabular}{@{}lll@{}}
\toprule
Logical demand & Left elevon & Right elevon\\
\midrule
Positive elevator only (\(a=0,e>0\)) & Up & Up\\
Positive aileron only (\(a>0,e=0\)) & Down & Up\\
Negative aileron only (\(a<0,e=0\)) & Up & Down\\
\bottomrule
\end{tabular}
\end{center}

Accordingly, positive aileron produces the differential pattern for a
right-roll command: left down, right up. Positive elevator produces equal
upward commands on both sides. These are mixer-level signs. The pilot's stick
direction depends on RC mapping and input configuration, while the final
physical movement depends on servo orientation, linkage, output type, and
\texttt{SERVO\textit{n}\_REVERSED}. The mixer comment describes the intended
non-reversed convention; it cannot guarantee that a particular airframe moves
correctly. Verify actual movement on the assembled aircraft before flight.

\section{Output scaling and configuration notes}

The mixer works on scaled logical values, then the servo-channel layer converts
them to PWM using each output's configured limits and type. For angle-style
outputs, positive scaled values increase PWM from trim unless the output is
reversed; reversal changes the final sign/direction. Mechanical travel and
endpoint configuration therefore affect the physical result, not the
logical mixing equation.

\texttt{MIXING\_GAIN} defaults to \(0.5\), reducing the chance that simultaneous
full roll and pitch inputs saturate an output; raising it increases travel but
can cause clipping. \texttt{MIXING\_OFFSET} changes the relative response to
roll and pitch as defined above. For unusual transmitter-side mixing,
\texttt{MANUAL\_RCMASK} can copy selected radio inputs directly to corresponding
outputs in MANUAL after the normal mixer runs. Those selected output channels
may consequently bypass the expected software elevon result in MANUAL.

\section{Source references}

\begin{itemize}
  \item \url{ArduPlane/mode_manual.cpp}: MANUAL input forwarding.
  \item \url{ArduPlane/Attitude.cpp}: stabilized roll and pitch outputs.
  \item \url{ArduPlane/servos.cpp}: elevon mixer, common output path, and
  \texttt{MANUAL\_RCMASK} handling.
  \item \url{ArduPlane/Parameters.cpp}: \texttt{MIXING\_GAIN},
  \texttt{MIXING\_OFFSET}, and \texttt{MANUAL\_RCMASK} parameter descriptions.
  \item \url{libraries/SRV\_Channel/SRV\_Channel.cpp}: scaled-output to PWM
  conversion and output reversal.
\end{itemize}

\end{document}
